% Generated draft: requires mathematical and source review.
\documentclass[11pt]{article}
\usepackage[margin=25mm]{geometry}
\usepackage{fontspec}
\setmainfont{Latin Modern Roman}
\usepackage{amsmath,amssymb,mathtools,unicode-math}
\setmathfont{Latin Modern Math}
\usepackage{enumitem,needspace,xurl,hyperref}
\setlength{\parindent}{0pt}
\setlength{\parskip}{5pt}
\setlength{\emergencystretch}{3em}
\newcommand{\N}{\mathbb{N}}
\newcommand{\Z}{\mathbb{Z}}
\newcommand{\Q}{\mathbb{Q}}
\newcommand{\R}{\mathbb{R}}
\newcommand{\C}{\mathbb{C}}
\newcommand{\F}{\mathbb{F}}
\newcommand{\A}{\mathbb{A}}
\newcommand{\PP}{\mathbb{P}}
\newcommand{\E}{\mathbb{E}}
\newcommand{\eps}{\varepsilon}
\newcommand{\dd}{\,\mathrm{d}}
\begin{document}
\section*{UnsolvedMath ID: 30005739}
\textbf{Source problem code:} catalogue entry
\subsection{Definitions and mathematical statement}
For terminology, let \(k\) be a field and let \(A\) be a \(k\)-vector space with a commutative associative bilinear product \(\cdot\). The source does not specify \(k\), its characteristic, whether \(A\) is finite-dimensional, or whether the product has an identity.

A Lie bracket on \(A\) is a bilinear alternating operation \([\ ,\ ]\) satisfying
\[
[x,[y,z]]+[y,[z,x]]+[z,[x,y]]=0
\qquad(x,y,z\in A).
\]
A family \(([\ ,\ ]_i)_{i\in I}\) of Lie brackets on the same vector space is called compatible if every finite linear combination
\[
[x,y]_{F,c}=\sum_{i\in F}c_i[x,y]_i,
\qquad F\subseteq I\text{ finite},\quad (c_i)_{i\in F}\in k^F,
\]
is a Lie bracket. This is a standard possible meaning of the source's reference to compatible Lie brackets; the supplied passage does not specify the intended family or its axioms.

For a single bracket, the Poisson compatibility identity is
\[
[x,y\cdot z]=[x,y]\cdot z+y\cdot[x,z]
\qquad(x,y,z\in A),
\]
whereas the transposed Poisson compatibility identity is
\[
2z\cdot[x,y]=[z\cdot x,y]+[x,z\cdot y]
\qquad(x,y,z\in A).
\]
These identities specify terminology for the draft; they do not supply the missing hypotheses of the research question. In particular, the characteristic of \(k\) matters for interpreting the factor \(2\).

\textbf{Source research problem.} Develop the theory requested in the source under the name ``transpose bi-Hamiltonian algebra'' (rendered ``transposed bi-Hamiltonian algebras'' in the registry). More generally, determine the algebraic structures obtained by replacing the single bracket of a Poisson algebra with compatible Lie brackets, with a bi-Hamiltonian system or a bi-Hamiltonian algebra intended as a special case. Ask the corresponding question for transposed Poisson algebras.

Under the compatible-family interpretation above, one possible working formulation is to study data
\[
\bigl(A,\cdot,([\ ,\ ]_i)_{i\in I}\bigr)
\]
for which each bracket satisfies the Poisson identity with the common product, or, in the transposed version, each bracket satisfies the transposed Poisson identity with that product. A two-bracket family is a possible starting point for interpreting the bi-Hamiltonian special case. However, the source does not establish that these axioms define its intended bi-Hamiltonian or transpose bi-Hamiltonian objects, nor whether further mixed identities or geometric data are required. The two compatibility identities above describe separate versions of the question and are not imposed simultaneously.

The supplied question is a research programme rather than a fully specified conjecture: no classification equivalence, universal property, proposed theorem, or precise meaning of ``develop the theory'' is given. Consequently, the working formulation is conditional and requires human mathematical review before it can be identified with the intended source problem.
\subsection{Short English statement}
The problem asks for a theory combining compatible Lie brackets with the multiplication of a Poisson or transposed Poisson algebra, including the intended bi-Hamiltonian special cases. The source leaves the underlying hypotheses, exact bi-Hamiltonian axioms, and desired classification or theorem unspecified, so the displayed framework is a conditional interpretation.
\subsection{Sources}
\begin{itemize}
\item Supplied local UnsolvedMath registry record, permanent ID \(30005739\), title ``Algebraic Structures for Transposed Bi-Hamiltonian Systems,'' collection ``Oberwolfach Reports Open Problems,'' source problem code \texttt{catalogue entry}. The registry records proposed year \(2023\); no individual proposer is supplied.
\item Supplied source citation: \emph{Mini-Workshop: Poisson and Poisson-type algebras} (2024), \emph{Oberwolfach Reports}, DOI \texttt{10.4171/owr/2023/46}, pp.\ 2681--2715. Supplied source URL and literature-source link, labelled ``Oberwolfach Report 46/2023'': \url{https://doi.org/10.4171/owr/2023/46}. The supplied background describes Oberwolfach Reports as records of mathematical workshops and their open-problem discussions. The citation's \(2024\) date and its \(2023\) report designation are retained without resolving their relationship.
\item The supplied original passage asks about transpose bi-Hamiltonian algebras, replacing the bracket in a Poisson algebra by compatible Lie algebra structures with bi-Hamiltonian systems or algebras as special cases, and making the analogous replacement for transposed Poisson algebras. No exact page within the supplied page range or individual attribution is provided.
\item Literature assessment dated exactly \texttt{2026-08-22}: the supplied curation labels the problem open and says that the report presents natural research questions and that the cited compatible-bracket literature does not supply a general theory of the transpose bi-Hamiltonian and transposed Poisson constructions. Its listed checked source is Oberwolfach Report 46/2023 at the DOI link above. This is supplied dated curation, not an independently verified literature assessment; the background expressly calls for independent verification of time-sensitive status and cited claims.
\item The registry labels the statement \texttt{corrected\_verified} and states that it was reconstructed and checked against the official report. These are supplied verification claims, not verification performed for this draft.
\end{itemize}
\end{document}
